\documentclass[10pt,a4paper]{article}
\usepackage[T1]{fontenc}
\usepackage[margin=2cm]{geometry}
\usepackage{isabelle,isabellesym}
\usepackage{amsmath,amssymb}
\usepackage{pdfsetup}

\urlstyle{rm}
\isabellestyle{it}

\begin{document}

\title{Dirac's Dense 6-Critical Graphs}
\author{Carles Mar\'in}
\maketitle

\begin{abstract}
A graph is $k$-critical if its chromatic number is $k$ and deleting any vertex or edge lowers it. We formalize Dirac's construction~\cite{Dirac1952}: the join of two odd cycles of length $2n+1$ is $6$-critical, has $4n+2$ vertices and $4n^2+8n+3$ edges, so the largest number $f_6(4n+2)$ of edges of a $6$-critical graph on $4n+2$ vertices is at least $4n^2+8n+3$. The statement follows the variant of Erd\H{o}s problem 917 in the formal-conjectures repository.
\end{abstract}

\paragraph{Acknowledgement.} The formalization and the text of this document were prepared with the
assistance of generative AI systems, and checked by the author.

\tableofcontents

\parindent 0pt\parskip 0.5ex

\input{session}

\bibliographystyle{abbrv}
\bibliography{root}

\end{document}
